\documentclass[12pt]{ctexart}
\usepackage{amsmath, amssymb, amsthm}
\usepackage{geometry}
\geometry{a4paper, margin=2.5cm}
\usepackage{enumitem}

\newtheorem{problem}{题目}
\newenvironment{solutionarea}{\framebox(500,80){}}{}

\title{线性代数 · 第一章 · 向量与向量空间 \\ 理论作业}
\author{}
\date{}

\begin{document}
\maketitle

\begin{center}
\textbf{说明}：在标注了 \texttt{\%\%\% 在此作答 \%\%\%} 的区域填写你的答案。
完成后将文件发给我批改。编程题部分见 \texttt{homework.py}。
\end{center}

% ──────────────────────────────────────────────
\begin{problem}[向量空间的判定 · ⭐]
判断以下集合是否为 $\mathbb{R}^2$ 的子空间。若是，给出它的维数和一组基；若不是，说明违反了什么条件。

(1) $W_1 = \{(x, y) \in \mathbb{R}^2 \mid y = -x\}$

(2) $W_2 = \{(x, y) \in \mathbb{R}^2 \mid x \geq 0\}$

(3) $W_3 = \{(x, y) \in \mathbb{R}^2 \mid y = x^2\}$
\end{problem}

%%% 在此作答 (1) %%%



%%% 在此作答 (2) %%%



%%% 在此作答 (3) %%%



% ──────────────────────────────────────────────
\begin{problem}[线性无关判定 · ⭐⭐]
判断以下向量组是否线性无关。请写出方程 $c_1\mathbf{v}_1 + c_2\mathbf{v}_2 + c_3\mathbf{v}_3 = \mathbf{0}$ 并求解，展示完整推导过程。

\[
\mathbf{v}_1 = \begin{bmatrix} 1 \\ 1 \\ 0 \end{bmatrix},\quad
\mathbf{v}_2 = \begin{bmatrix} 0 \\ 1 \\ 1 \end{bmatrix},\quad
\mathbf{v}_3 = \begin{bmatrix} 1 \\ 0 \\ -1 \end{bmatrix}
\]

\end{problem}

%%% 在此作答 %%%



% ──────────────────────────────────────────────
\begin{problem}[求基与维数 · ⭐⭐]
求以下 $\mathbb{R}^4$ 的子空间 $W$ 的一组基和维数：

\[
W = \{(x_1, x_2, x_3, x_4) \in \mathbb{R}^4 \mid x_1 + x_2 + x_3 + x_4 = 0\}
\]

（提示：用自由变量表示约束变量，提取基向量。）
\end{problem}

%%% 在此作答 %%%



% ──────────────────────────────────────────────
\begin{problem}[线性组合与 Span · ⭐⭐]
设 $\mathbf{v}_1 = [1, 2, 1]^T$, $\mathbf{v}_2 = [2, 5, 0]^T$, $\mathbf{v}_3 = [3, 3, 8]^T$。

问 $\mathbf{w} = [1, 1, 1]^T$ 是否属于 $\operatorname{span}\{\mathbf{v}_1, \mathbf{v}_2, \mathbf{v}_3\}$？
如果是，把它写成 $\mathbf{v}_1, \mathbf{v}_2, \mathbf{v}_3$ 的线性组合。
\end{problem}

%%% 在此作答 %%%



% ──────────────────────────────────────────────
\begin{problem}[余弦相似度的几何解释 · ⭐⭐]
(1) 设 $\mathbf{u} = [1, 2]^T$, $\mathbf{v} = [4, -2]^T$。计算它们的内积、范数、夹角余弦、夹角（度数，保留 1 位小数）。

(2) 在机器学习中，为什么用余弦相似度而不是欧氏距离来衡量文本向量的相似性？用一句话回答。

\end{problem}

%%% 在此作答 (1) %%%



%%% 在此作答 (2) %%%



% ──────────────────────────────────────────────
\begin{problem}[子空间交 · ⭐⭐]
设 $U = \operatorname{span}\{[1, 1, 0]^T, [0, 1, 1]^T\}$, $W = \operatorname{span}\{[1, 0, 1]^T, [0, 1, -1]^T\}$。

求 $U \cap W$ 的一组基和维数。（提示：$U \cap W$ 中的向量必须同时属于两个 span。）
\end{problem}

%%% 在此作答 %%%



% ──────────────────────────────────────────────
\begin{problem}[范数不等式 · ⭐⭐⭐]
证明以下不等式对任意 $\mathbf{u}, \mathbf{v} \in \mathbb{R}^n$ 成立：

(1) Cauchy-Schwarz 不等式：$|\mathbf{u} \cdot \mathbf{v}| \leq \|\mathbf{u}\| \|\mathbf{v}\|$

（提示：考虑 $\|\mathbf{u} - t\mathbf{v}\|^2 \geq 0$ 对任意 $t \in \mathbb{R}$ 成立。）

(2) 三角不等式：$\|\mathbf{u} + \mathbf{v}\| \leq \|\mathbf{u}\| + \|\mathbf{v}\|$

（用 (1) 的结论。）
\end{problem}

%%% 在此作答 (1) %%%



%%% 在此作答 (2) %%%



% ──────────────────────────────────────────────
\begin{problem}[线性无关与基的扩张 · ⭐⭐⭐]
设 $\mathbf{v}_1, \mathbf{v}_2 \in \mathbb{R}^3$ 线性无关，且 $\mathbf{v}_3 \notin \operatorname{span}\{\mathbf{v}_1, \mathbf{v}_2\}$。

(1) 证明 $\{\mathbf{v}_1, \mathbf{v}_2, \mathbf{v}_3\}$ 是 $\mathbb{R}^3$ 的一组基。

(2) 给定 $\mathbf{v}_1 = [1, 0, -1]^T$, $\mathbf{v}_2 = [2, 1, 0]^T$，找一个 $\mathbf{v}_3$ 使得 $\{\mathbf{v}_1, \mathbf{v}_2, \mathbf{v}_3\}$ 是 $\mathbb{R}^3$ 的基，并验证。
\end{problem}

%%% 在此作答 (1) %%%



%%% 在此作答 (2) %%%



\end{document}
